\documentclass[../../main.tex]{subfiles}
\begin{document}

{\bf [解答]}

\paragraph{(1)}

$f(x)=0$ は実数範囲で少なくとも1つ解を持つことを示せば良い．
$y=f(x)$ のグラフを考えると無限遠で$f(x) \to \pm \infty \ (x \to \pm \infty)$である．これと$f(x)$ は連続であることから，中間値の定理から $f(x)=0$ は実数範囲で少なくとも1つ解$x=-A$を持つ．従って因数定理によって因数分解できて
\begin{align}
  f(x) = (x+A)(x^2+Bx+C)
\end{align}
と書ける．よって題意は示された．

\paragraph{(2)}

$f(x)$の係数は全て実数だから，複素数 $\alpha$ に対して $f(\alpha)=0$ ならば $f(\bar{\alpha})=0$ である．
$f(x)=0$ は三次式なので複素数範囲で重根込めて3つの解しかないことから $f(x)=0$ の解は以下のいずれかである．

\begin{enumerate}
  \item[$1^\circ$] 3つとも実解
  \item[$2^\circ$] 1つ実解，2つが複素解
\end{enumerate}

それぞれの場合について題意を示す．

\subparagraph{$1^\circ$ 3つとも実解の時}

3つの解 $\alpha, \beta, \gamma$ ($\alpha, \beta, \gamma < 0$) とすると，
因数定理から
\begin{align}
  f(x)=(x-\alpha)(x-\beta)(x-\gamma)
\end{align}
とおけるから，係数比較して
\begin{align}
  \begin{cases}
    a = -(\alpha+\beta+\gamma) > 0 \\
    b = \alpha\beta+\beta\gamma+\gamma\alpha  > 0 \\
    c = -\alpha\beta\gamma  > 0
  \end{cases}
\end{align}
である．よって，$a, b, c > 0$である．

\subparagraph{$2^\circ$ 1つ実解，2つが複素解の時}

実数解を $\alpha$，共役解を $p \pm qi$ ($p, q \in \mathbb{R}$, $p < 0, \alpha < 0$) とおくと，$f(x)$は
\begin{align}
  f(x) &= (x-\alpha)(x-(p+qi))(x-(p-qi)) \\
  &= (x-\alpha)(x^2 - 2px + (p^2+q^2))
\end{align}
と表せる．係数比較して
\begin{align}
  \begin{cases}
    a = -\alpha - 2p  > 0 \\
    b = p^2+q^2+2\alpha p  > 0 \\
    c = -\alpha(p^2+q^2)  > 0
  \end{cases}
\end{align}
だから $a, b, c > 0$となる．
\casesend

以上から，いずれの場合も $a, b, c > 0$ である．

{\bf [別解]}

\paragraph{(1)}

代数学の基本定理より，方程式 $f(x)=0$ は複素数範囲で3つの解をもつ．
$f(x)$ は実数係数多項式であるから，複素数 $\alpha$ が $f(\alpha)=0$ をみたせば，その共役複素数 $\bar{\alpha}$ も $f(\bar{\alpha})=0$ をみたす．
したがって虚数解が存在するならば必ず共役なペアで現れるため，虚数解の個数は偶数（0個または2個）である．
ゆえに，3つの解のうち少なくとも1つは実数解でなければならない．
その実数解の1つを $\alpha$ とおくと，因数定理より $f(x)$ は実数係数の1次式 $(x-\alpha)$ で割り切れる．
実数係数の多項式同士の割り算であるから，商 $x^2+Bx+C$ の係数 $B, C$ もすべて実数である．
したがって $A = -\alpha$ とおけば，
\begin{align}
  f(x) = (x+A)(x^2+Bx+C)
\end{align}
となる実数 $A, B, C$ が存在する．

\newpage

\end{document}
